% SPDX-License-Identifier: Apache-2.0
\section{Supply Chain: A Schema-Validated CBOM}
\label{sec:cbom}

\subsection{Why a CBOM, and why the standard matters}

A Cryptographic Bill of Materials is the natural artefact for the question
``which cryptographic assets are in this system''. CycloneDX 1.6~\cite{cyclonedx16} extends the
SBOM format with \texttt{cryptoProperties}, giving components a place to declare
their algorithm properties, their certification status, and their occurrences.

The temptation is to invent a format. It is much easier, because a
project-specific JSON schema can be designed to fit exactly what the generator
already emits, and then validation is guaranteed to pass. That is not
validation; it is transcription.

We use the official CycloneDX 1.6 schema~\cite{cyclonedx-schema16}, vendored into
the repository with its hash recorded, and validate against it in CI on every push.

\subsection{The generator}

\texttt{scripts/build-cbom.js} reads a scan of the source tree
(\texttt{tools/crypto-audit-scan.json}, 370 entries) and emits
\texttt{tools/cbom-cyclonedx.json}. It is deterministic: the serial number is
derived by hashing the component set and the root properties, and the timestamp
honours \texttt{SOURCE\_DATE\_EPOCH}, so two runs over the same tree produce
byte-identical output. A \texttt{-{}-check} mode regenerates and compares, and
fails if the committed artefact is stale.

\subsection{What is a component, and what is not}

The scan finds more than cryptographic primitives. It also finds a formal
specification language (TLA+), a mathematical transform (the number-theoretic
transform), and a protocol construction (Double Ratchet). None of these is a
cryptographic asset in the sense the schema means, and emitting them as
components would inflate the inventory with things no one can migrate.

They are not silently dropped. They are recorded as root properties in the
\texttt{fibemate:} namespace with the prefix \texttt{excludedAsset}, each with
the reason for exclusion attached, so that the decision is auditable and a
reader can disagree with it.

The result is 9 components:

\begin{center}
\begin{tabular}{@{}llrl@{}}
\toprule
Component & Primitive & Occurrences & Quantum status \\
\midrule
ML-KEM & kem & 164 & safe \\
SHA-256 & hash & 94 & safe \\
SM2 & signature & 90 & vulnerable \\
P-256/ECDH & key-agree & 72 & vulnerable \\
SM3 & hash & 62 & safe \\
ML-DSA/fml-dsa & signature & 53 & safe \\
SM4 & block-cipher & 40 & safe \\
AES & block-cipher & 39 & safe \\
SLH-DSA & signature & 33 & safe \\
\bottomrule
\end{tabular}
\end{center}

Across 304 distinct files, with 647 component--file occurrences.

\subsection{Two numbers that must not be confused}

647 is the number of occurrences. 304 is the number of distinct files. They
differ because 139 files mention more than one algorithm, and one HTML page
mentions nine.

This distinction is not pedantry. Our viewer originally displayed 647 under a
label reading ``total files'', overstating the footprint by 343. The error
survived because the number was hardcoded and therefore never recomputed, and
because no one checks a number that looks plausible.

\subsection{What the schema forced us to stop saying}

The most useful thing about validating against an external schema is that it
refuses claims the project had been making.

Our first CBOM failed validation with 51 errors in 7 categories. They were not
typos. They were category errors: risk information placed where the schema does
not allow it, an \texttt{implementationPlatform} value (``javascript'') that is
not a platform, a \texttt{certificationLevel} of
\texttt{known-answer-tested} which is not a certification level the schema
recognises, and primitive names in upper case where the enumeration is lower
case.

Each of these was a small overclaim. ``KAT passed'' is not a certification.
JavaScript is not a hardware platform. Fixing them meant either finding the
correct field or \emph{removing the claim}, and in several cases the honest
answer was removal: the project holds no FIPS or CC certification, so the
certification level is \texttt{none}, and it says so.

The design rule we settled on: \textbf{a missing field is better than an
invented field}. A gap is visible and can be closed. A plausible fabrication is
invisible and propagates.

\subsection{Continuous validation}

The CBOM is validated on every push in two ways: against the vendored official
schema, and against a drift check that regenerates the artefact and fails if the
committed copy differs. The drift check also verifies that the version recorded
in the CBOM matches \texttt{package.json}, because a bill of materials that
misstates the version of the thing it describes is worse than no bill of
materials.
